\documentclass{article}
\usepackage{amsmath}
\usepackage{cancel}
\usepackage{graphicx} % for creating figures

\title{The binary model system}

\begin{document}
\maketitle
Consider a system of $N$ quantum particles which each may have one of two
possible energies denoted $+\varepsilon$ and $-\varepsilon$. This model applies
reasonably well to paramagnetic materials which consist of a collection of
dipoles that we approximate as having magnetic moments $+m$ and $-m$ which
correspond to pointing either \textit{either up or down}. A state in this
system can be easily represented by a graphic like: 
\begin{equation}
  \uparrow \downarrow \uparrow \uparrow \uparrow \downarrow \ldots
\end{equation}
where each arrow represents the state for one of the particles in the system.

If there are $N$ particles, then there are $2^N$ possible states with each
distinct state given as a term in
\begin{equation}
  \Pi_i (\uparrow_i + \downarrow_i)
\end{equation}
For example, if $N=2$ then there are $2^2=4$ possible states. They are given by
the terms of
\begin{equation}
  (\uparrow_1 + \downarrow_1)(\uparrow_2+\downarrow_2) = \uparrow_1\uparrow_2 + \uparrow_1\downarrow_2 + \downarrow_1\uparrow_2 + \downarrow_1\downarrow_2
\end{equation}
The product $\Pi_i(\uparrow_i+\downarrow_i)$ is called the \textit{generating
  function} as it generates all of the possible states of the system.

If we consider the entire collection of dipoles, then we can define the total
magnetic moment $M$ which is just the sum of all of the magnetic moments of the
system. Clearly, $M$ must take a value between $-Nm$ and $Nm$ in integer steps.
This indicates that there are $N+1$ possible values of the total magnetic moment
(don't forget to count 0!). For systems with large values of $N$ we can clearly
see that there are many more possible states than there are total magnetic
moments.

For a given energy of the system, it will prove useful to find the analytic form of the
multiplicity. We shall denote the \textit{multiplicity function} as
\begin{equation}
  g(N, s)
\end{equation}
where $N$ is the total number of particles in the system, and $s$ is the index
specifying the energy. For the system of $N$ binary magnetic moments, the energy
of the collection is independent of the order in which we specify the values of
$m$ for each particle e.g. $E(\uparrow \downarrow) = E(\downarrow \uparrow)$.
For simplicity, let's assume $N$ to be a large, even number and define the
difference between the number of particles in the $\uparrow$ state and the
number of particles in the $\downarrow$ state to be $2s$: 
\begin{equation}
  N_{\uparrow} - N_{\downarrow} = 2s
\end{equation}
where the $2$ is just for convenience. Then, we can say that
the coefficients in the expansion
\begin{equation}
  (\uparrow  +\downarrow)^N
\end{equation}
represent the multiplicity of each distinct energy state. This is because we are
no longer tracking the ordering, and therefore: 
\begin{equation}
  (\uparrow+\downarrow)^2 = \uparrow\uparrow + 2\uparrow\downarrow + \downarrow\downarrow.
\end{equation}
From this equation, we would say that $g(2, 1) = 1$, $g(2, 0) = 2$, and
$g(2,-1)=1$. For a general $N$, the coefficients can be found as the
coefficients of the binomial expansion (i.e. the elements of Pascal's triangle).
Since we must also have that $N = N_{\uparrow}+N_{\downarrow}$, we can say
that
\begin{equation}
  N_{\uparrow} = \frac{1}{2}N+s \qquad N_{\downarrow} = \frac{1}{2}N-s. 
\end{equation}
The number of particles in the $\uparrow$ state can range from $N_\uparrow = 0$
to $N_\uparrow = N$, and therefore, $s$ takes on values $s=-\frac{1}{2}N$ to
$s=\frac{1}{2}N$. 
\begin{align}
  (\uparrow+\downarrow)^N &= \sum\limits_{s=-\frac{1}{2}N}^{\frac{1}{2}N}\frac{N!}{N_{\uparrow}!(N-N_{\uparrow})!}\uparrow^{N_\uparrow}\downarrow^{N-N_{\uparrow}} \\
                          &=\sum\limits_{s=-\frac{1}{2}N}^{\frac{1}{2}N}\frac{N!}{(\frac{1}{2}N+s)!(\frac{1}{2}N-s)!}\uparrow^{\frac{1}{2}N+s}\downarrow^{\frac{1}{2}N-s} 
\end{align}
Identifying the multiplicity function as the coefficients in the above equation,
we conclude that
\begin{equation}
  g(N,s) = \frac{N!}{(\frac{1}{2}N+s)!(\frac{1}{2}N-s)!} = \frac{N!}{N_{\uparrow}!N_{\downarrow}!}
\end{equation}


\end{document}




